\chapter{Actions and Recipes: Functions, Closures, and Pipelines}

\section{What is a Function?}

\begin{kidbox}{Mom's Secret Pancake Recipe}
Imagine you love pancakes. Every Saturday morning, you want delicious, fluffy blueberry pancakes!
Do you sit down with a blank piece of paper and invent how to make pancakes from scratch every single time?
No! You open mom's cookbook to the page titled \textbf{"make\_pancakes"}.
The recipe says:
\begin{enumerate}
    \item Give me ingredients: 2 eggs, 1 cup flour, half cup milk.
    \item Mix them in a bowl, pour onto hot pan, flip when bubbly.
    \item Hand back a warm plate of golden pancakes!
\end{enumerate}
In Shae, a \textbf{Function} is a named recipe that you can use whenever you want:
\begin{lstlisting}[language=shaelang]
fn make_pancakes(eggs, cups_flour, cups_milk) {
    let batter = eggs + cups_flour + cups_milk;
    return "Plate of " + batter + " delicious pancakes!";
}

let breakfast = make_pancakes(2, 1, 1);
println(breakfast);
\end{lstlisting}
\end{kidbox}

\section{First-Class Functions: Passing Recipes to Friends}

In Shae, functions are \textbf{first-class citizens}. This means functions can be stored in variables, passed into other functions as arguments, and returned from functions just like integers or strings!

\begin{lstlisting}[language=shaelang]
fn double(x) { x * 2 }
fn square(x) { x * x }

fn apply_math(n, math_recipe) {
    return math_recipe(n);
}

println(apply_math(5, double)); // 10
println(apply_math(5, square)); // 25
\end{lstlisting}

\section{Closures: The Function's Memory Backpack}

\begin{kidbox}{The Little Explorer's Backpack}
Imagine a little explorer named Finn who is about to go on a hike up Mount Everest.
Before leaving home, his mother gives him a warm thermos of hot cocoa and puts it inside his backpack.
No matter how high Finn climbs into the snowy mountains, whenever he stops to take a break, he can open his backpack and sip the warm cocoa from home!
In Shae, a \textbf{Closure} is a function that carries a backpack full of variables from the place where it was born!
\begin{lstlisting}[language=shaelang]
fn make_counter(start_value) {
    let count = start_value;
    
    // This inner function packs 'count' into its backpack!
    return fn() {
        count = count + 1;
        return count;
    };
}

let counter_a = make_counter(10);
let counter_b = make_counter(100);

println(counter_a()); // 11
println(counter_a()); // 12
println(counter_b()); // 101
\end{lstlisting}
\end{kidbox}

\begin{deepbox}{Upvalues and Closed Enclosing Frames}
In systems compilers, lexical closures present a fundamental challenge: when the outer function \texttt{make\_counter} returns, its stack frame is popped. If the inner closure references \texttt{count}, accessing the stack address would cause an undefined memory read (dangling stack pointer).

Shae resolves this via \textbf{Upvalues}:
\begin{itemize}
    \item \textbf{Open Upvalue}: While the outer function is executing, the upvalue points directly to the active stack slot (\texttt{\&mut vm.stack[slot]}).
    \item \textbf{Closed Upvalue}: When the outer stack frame pops, the VM migrates the captured variable from the stack into heap-allocated reference-counted memory, updating the upvalue pointer to point to the heap copy.
\end{itemize}
This ensures zero allocation overhead for non-escaping variables, while guaranteeing complete memory safety for long-lived closures.
\end{deepbox}

\section{The Pipeline Operator: The Water Slide of Data}

\begin{kidbox}{The Water Park Funnel}
Have you ever been to a water park?
You don't take a bucket of water and walk it up the stairs one step at a time.
You jump into the water slide!
First you zoom through the blue tunnel, then you slide through the giant funnel, and finally you splash into the swimming pool!
In Shae, the \textbf{Pipeline Operator} (\texttt{|>}) lets your data slide smoothly through your functions:
\begin{lstlisting}[language=shaelang]
let numbers = [1, 2, 3, 4, 5, 6, 7, 8, 9, 10];

let result = numbers
    |> filter(fn(x) { x % 2 == 0 })  // Only evens: [2, 4, 6, 8, 10]
    |> map(fn(x) { x * x })          // Square them: [4, 16, 36, 64, 100]
    |> sum();                        // Sum: 220

println("Final sum: " + result);
\end{lstlisting}
Without pipelines, you would have to write inside-out math: \texttt{sum(map(filter(numbers, ...), ...))}.
With pipelines, it reads like a clear, flowing story from left to right!
\end{kidbox}
